\documentclass[conference]{IEEEtran}
\IEEEoverridecommandlockouts
\usepackage{cite}
\usepackage{amsmath,amssymb,amsfonts}
\usepackage{algorithmic}
\usepackage{graphicx}
\usepackage{textcomp}
\usepackage{xcolor}
\usepackage{fancyhdr}

\pagestyle{fancy}
\fancyhf{}
\renewcommand{\headrulewidth}{0pt}
\fancyfoot[L]{\textit{ERU Symposium 2026 - DOI (to be filled by organisers)}}
\renewcommand{\footrulewidth}{0pt}

\def\BibTeX{{\rm B\kern-.05em{\sc i\kern-.025em b}\kern-.08em
    T\kern-.1667em\lower.7ex\hbox{E}\kern-.125emX}}

\begin{document}

\title{Intelligent Humanoid Airport Concierge: A RAG-Based Voice Interface and Autonomous Navigation System for the Unitree G1}

\author{
\IEEEauthorblockN{M.A.A. Rafi}
\IEEEauthorblockA{
\textit{Department of Computer Science and Engineering} \\
\textit{University of Moratuwa} \\
Katubedda, Sri Lanka \\
Email: abdulr.23@cse.mrt.ac.lk}
\and
\IEEEauthorblockN{Sulochana Sooriyaarachchi}
\IEEEauthorblockA{
\textit{Department of Computer Science \& Engineering} \\
\textit{University of Moratuwa} \\
Katubedda, Sri Lanka \\
Email: sulochanas@cse.mrt.ac.lk}
}

\maketitle

\begin{IEEEkeywords}
humanoid robotics, retrieval-augmented generation, autonomous navigation, LiDAR SLAM, multilingual speech interface
\end{IEEEkeywords}

\section{Introduction}
Airport terminals need both natural-language guidance and physical wayfinding, a combination existing systems only partly cover: chatbots answer questions but cannot escort a passenger, while navigation robots move reliably but understand little natural language. This work integrates both on a Unitree G1 EDU humanoid. A retrieval-augmented voice interface answers passenger queries from an approved knowledge base, and a LiDAR-based navigation stack then guides the passenger to the identified destination. The system targets English, Sinhala and Tamil and is validated in MuJoCo/ROS 2 simulation.

\section{Literature Review}
Retrieval-augmented generation (RAG) grounds language-model outputs in an external document store, reducing hallucination and allowing knowledge updates without retraining \cite{lewis2020rag}. The ROS 2 Nav2 stack offers a modular, behaviour-tree-driven planner and controller demonstrated on service robots operating alongside pedestrians \cite{macenski2020marathon2}. Multilingual transcription remains a bottleneck for lower-resource languages such as Sinhala and Tamil; Whisper generalises well but can trail cloud recognisers on noisy, accented speech \cite{radford2022whisper}, and faster-whisper \cite{systran_fasterwhisper} reimplements it to cut inference time and memory. Few prior systems combine RAG conversation, multilingual speech and full-body humanoid navigation in one airport-concierge pipeline.

\section{Materials and Methods}

\subsection{Problem Formulation and Objectives}
The target scenario is a passenger asking, for example, ``Where is the taxi pickup point?'' The robot must transcribe the request, retrieve approved information, give a grounded spoken answer, identify the destination on its map, and guide the passenger there autonomously. The objectives are to: (1) build a voice interface for the G1; (2) ground answers in approved documents; (3) support English, Sinhala and Tamil; (4) map, localize and navigate to a selected destination, validated in simulation first; and (5) evaluate accuracy, latency, navigation success and computational cost.

\subsection{System Architecture}
Fig.~\ref{fig:architecture} shows the two pipelines. In the conversation pipeline, WebRTC voice-activity detection (VAD) segments speech, which is transcribed and passed to a LangChain agent. The agent retrieves context from a FAISS index of approved documents (Gemini embeddings) and a Gemini model generates a grounded spoken answer. If the request names a destination, the agent forwards it as a Nav2 goal, planned on a LiDAR-built map, localized through the frame chain odom $\rightarrow$ pelvis $\rightarrow$ mid360\_link, and executed as velocity commands on the simulated G1.

\begin{figure}[htbp]
\centering
% \includegraphics[width=\columnwidth]{system_architecture.png}
\caption{System architecture: the conversation pipeline (top) passes a spoken destination to the ROS 2 localization and navigation pipeline (bottom).}
\label{fig:architecture}
\end{figure}

\subsection{Conversational RAG Pipeline}
Table~\ref{tab:methods} lists the components chosen for the conversation pipeline; speech recognition uses Whisper, with faster-whisper \cite{systran_fasterwhisper} as a lighter local backend. RAG was preferred over fine-tuning because the knowledge base can be updated by re-indexing documents without retraining, which matters where airport signage, routes and services change.

\begin{table}[htbp]
\caption{Method Selection and Justification}
\label{tab:methods}
\centering
\begin{tabular}{|p{1.6cm}|p{2.1cm}|p{3.0cm}|}
\hline
\textbf{Component} & \textbf{Method} & \textbf{Justification} \\
\hline
Speech detection & WebRTC VAD & Lightweight; avoids processing silence \\
\hline
Speech recognition & Whisper / faster-whisper; cloud STT & Multilingual (EN/SI/TA) with a local backup \\
\hline
Retrieval and grounding & Gemini embeddings + FAISS; RAG with Gemini & Efficient search; restricts answers to approved sources \\
\hline
Orchestration & LangChain & Connects conversation, retrieval and navigation \\
\hline
Mapping / navigation & LiDAR SLAM + Nav2 & Mature, modular indoor navigation framework \\
\hline
\end{tabular}
\end{table}

\subsection{Localization and Navigation}
Indoor maps are built in MuJoCo/ROS 2 simulation from a simulated Livox MID-360 LiDAR, an Intel RealSense D435 depth camera, an IMU and ground-truth odometry (Fig.~\ref{fig:g1sim}). Given a 2D map and an initial pose, Nav2 plans a global path, converts it into local velocity commands and continuously updates the pose and obstacle estimate. To keep the humanoid stable, velocity limits are set conservatively at about 0.2 m/s linear and 0.3 rad/s angular.

\begin{figure}[htbp]
\centering
% \includegraphics[width=\columnwidth]{g1_simulation.png}
\caption{(a) Unitree G1 EDU humanoid \cite{unitree_g1}; (b) Nav2 costmap and LiDAR scan in RViz; (c) the G1 walking in the MuJoCo simulation.}
\label{fig:g1sim}
\end{figure}

\section{Results and Discussion}
The G1 MuJoCo simulation, simulated LiDAR and depth sensing, ROS 2 odometry/TF, voxel mapping, RViz visualization and text-to-speech output are complete. Autonomous navigation to a selected goal and English RAG conversation have been demonstrated end-to-end, with grounded English answers and reliable mapping and goal navigation in simulation. Sinhala and Tamil recognition is still being tuned: local Whisper models are less reliable for these languages, and a cloud recogniser improves them. Running perception and simulation together slows navigation, and 3D obstacle avoidance with the depth camera is not yet implemented, so the 3D LiDAR may miss certain types of obstacles.

Fig.~\ref{fig:benchmark} quantifies the speech backend from the published faster-whisper benchmark \cite{systran_fasterwhisper}, as relative reduction $\eta = (X_b - X_o)/X_b \times 100\%$ and speedup $S = T_b/T_o$ ($b$: baseline, $o$: optimized). On CPU (small model), faster-whisper INT8 cuts 13~min of transcription from 418~s to 102~s (75.6\% lower, $S = 4.10\times$) and RAM from 2335 to 1477~MB (36.7\% lower). On GPU (large-v2), time falls from 143~s to 59~s (58.7\% lower, $S = 2.42\times$) and VRAM from 4708 to 2926~MB (37.9\% lower). These are component-level values, not end-to-end G1 measurements; whole-system latency and power draw remain to be measured on the robot.

\begin{figure}[htbp]
\centering
% \includegraphics[width=\columnwidth]{whisper_benchmark.png}
\caption{faster-whisper (INT8) vs. OpenAI Whisper, from the published benchmark \cite{systran_fasterwhisper}. Component-level, not end-to-end G1 results.}
\label{fig:benchmark}
\end{figure}

\section{Conclusion and Future Work}
A single humanoid platform can combine grounded conversational AI with autonomous indoor navigation for an airport-concierge role, with English conversation and simulated goal navigation working end-to-end. Future work will link spoken destinations to Nav2 goals, replace ground-truth odometry with LiDAR-inertial localization, improve Sinhala and Tamil recognition, add depth-based 3D obstacle detection, and move to the physical G1 for tests with real users.

\begin{thebibliography}{00}
\bibitem{lewis2020rag} P. Lewis et al., ``Retrieval-augmented generation for knowledge-intensive NLP tasks,'' in \textit{Adv. Neural Inf. Process. Syst. (NeurIPS)}, vol. 33, 2020, pp. 9459--9474.
\bibitem{macenski2020marathon2} S. Macenski, F. Mart\'in, R. White, and J. Gin\'es Clavero, ``The Marathon 2: A navigation system,'' in \textit{Proc. IEEE/RSJ Int. Conf. Intelligent Robots and Systems (IROS)}, 2020.
\bibitem{radford2022whisper} A. Radford, J. W. Kim, T. Xu, G. Brockman, C. McLeavey, and I. Sutskever, ``Robust speech recognition via large-scale weak supervision,'' arXiv:2212.04356, 2022.
\bibitem{unitree_g1} Unitree Robotics, ``Unitree G1 humanoid robot,'' 2024. [Online]. Available: unitree.com/g1
\bibitem{systran_fasterwhisper} SYSTRAN, ``faster-whisper: Faster Whisper transcription with CTranslate2,'' GitHub. [Online]. Available: github.com/SYSTRAN/faster-whisper
\end{thebibliography}

\end{document}
